% ==========================================================================
%  Voxel Royale Distribuido - Rascunho do relatorio da Entrega 1
% --------------------------------------------------------------------------
%  Para a SUBMISSAO FINAL, troque a classe abaixo pelo template oficial da SBC
%  (sbc-template): baixe "sbc-template.sty" do modelo de artigos da SBC
%  (https://www.sbc.org.br) e use \documentclass[12pt]{article} +
%  \usepackage{sbc-template}. Este rascunho usa apenas a classe article para
%  compilar de forma portatil em qualquer instalacao LaTeX, mantendo a mesma
%  estrutura de secoes exigida pela Entrega 1 (problema, arquitetura, requisitos
%  implementados, detalhes de implementacao, desafios e papeis dos alunos).
% ==========================================================================
\documentclass[12pt]{article}

\usepackage[utf8]{inputenc}
\usepackage[T1]{fontenc}
\usepackage{lmodern}
\usepackage[brazil]{babel}
\usepackage[a4paper,margin=2.8cm]{geometry}
\usepackage{enumitem}
\usepackage{hyperref}
\hypersetup{hidelinks}

\setlength{\parskip}{0.4em}
\setlength{\parindent}{0pt}

\title{\textbf{Voxel Royale Distribuído:\\ um esqueleto de \textit{battle royale} em tempo real\\ com gRPC e \textit{web services}}}

\author{
Grupo de Computação Distribuída --- PUC Minas\\
\small Aluno 1 \texttt{PLACEHOLDER}, Aluno 2 \texttt{PLACEHOLDER}, Aluno 3 \texttt{PLACEHOLDER},\\
\small Aluno 4 \texttt{PLACEHOLDER}, Aluno 5 \texttt{PLACEHOLDER}, Aluno 6 \texttt{PLACEHOLDER},\\
\small Aluno 7 \texttt{PLACEHOLDER}, Aluno 8 \texttt{PLACEHOLDER}, Aluno 9 \texttt{PLACEHOLDER}
}
\date{}

\begin{document}
\maketitle

\begin{abstract}
\noindent\textbf{Resumo.} Este relatório descreve a primeira entrega do Voxel Royale
Distribuído, um mini \textit{battle royale} 2D projetado para rodar no navegador do
celular. O foco da Entrega 1 é demonstrar uma arquitetura distribuída clara que satisfaça,
obrigatoriamente, dois requisitos do enunciado: comunicação interna via \textbf{gRPC/RPC} e
exposição de \textbf{\textit{web services}} HTTP. Apresentamos a divisão em três
microsserviços Go --- Gateway, Lobby e Game ---, os contratos \textit{Protocol Buffers}, o
fluxo de mensagens e o empacotamento reprodutível em Docker Compose.

\vspace{0.6em}
\noindent\textbf{Abstract.} This report describes the first delivery of Voxel Royale
Distribuído, a mobile-first 2D battle royale. The goal of this delivery is to
demonstrate a clear distributed architecture satisfying two mandatory course requirements:
internal \textbf{gRPC/RPC} communication and HTTP \textbf{web services}. We present three Go
microservices (Gateway, Lobby and Game), the Protocol Buffers contracts, the message flow and
a reproducible Docker Compose packaging.
\end{abstract}

\section{Problema}

O objetivo do trabalho é construir, de forma jogável e mensurável, um sistema distribuído em
tempo real no qual até 50 jogadores participam de uma partida \textit{battle royale} voxel. Os
jogadores entram por QR Code, informam apenas um nome (sem conta ou autenticação) e disputam
partidas rápidas de até cinco minutos. O desafio central, do ponto de vista de Computação
Distribuída, não é o jogo em si, mas evidenciar conceitos da disciplina: comunicação entre
processos, contratos tipados, concorrência, um servidor autoritativo sobre o estado e um
caminho claro para observabilidade, tolerância a falhas e teste de carga nas próximas fases.

A Entrega 1 exige a implementação de pelo menos dois requisitos do enunciado. O grupo escolheu
\textbf{gRPC/RPC} para a comunicação interna entre serviços e \textbf{\textit{web services}}
HTTP para a operação e a entrada de jogadores. Esta entrega entrega um \textit{esqueleto
distribuído} executável que prova esses dois requisitos ponta a ponta.

\section{Arquitetura}

O sistema é um monorepo Go dividido em três microsserviços que se comunicam por gRPC e são
orquestrados por Docker Compose. O \textbf{Gateway} é a borda HTTP pública; o \textbf{Lobby}
gerencia o ciclo de vida das salas; e o \textbf{Game} é o servidor autoritativo da partida.

\begin{verbatim}
  Navegador / curl
        | HTTP JSON
        v
    Gateway :8080  --- gRPC --->  Lobby :50052   (salas)
        |
        +--------- gRPC --------> Game  :50051   (partida autoritativa)
\end{verbatim}

O Gateway traduz requisições HTTP/JSON em chamadas gRPC usando \textit{grpc-gateway}: as
anotações \texttt{google.api.http} nos arquivos \texttt{.proto} definem cada rota REST a partir
de um método RPC. Assim, o mesmo contrato \textit{Protocol Buffers} serve tanto para a
comunicação interna tipada (gRPC) quanto para a API pública (web service), evitando divergência
entre as duas interfaces. Os serviços se descobrem por DNS interno do Compose
(\texttt{game:50051}, \texttt{lobby:50052}) e cada um expõe um \texttt{/healthz} usado pelos
\textit{healthchecks}.

\section{Requisitos implementados}

\textbf{gRPC/RPC (requisito 1).} Os contratos vivem em \texttt{proto/lobby/v1/lobby.proto} e
\texttt{proto/match/v1/match.proto}, dos quais geramos o código Go com \texttt{protoc}. O Lobby
implementa o serviço \texttt{LobbyService} com seis RPCs: \texttt{CreateRoom}, \texttt{JoinRoom},
\texttt{GetRoom}, \texttt{StartRoom}, \texttt{LeaveRoom} e \texttt{SetReady}. O Game implementa
\texttt{GameService.StreamMatch}, que recebe o \textit{input} do jogador e devolve um
\textit{snapshot} autoritativo da partida (posições, vida, armas, baús, zona segura e ranking).

\textbf{Web services HTTP (requisito 2).} O Gateway publica, na porta 8080, web services REST/JSON
mapeados diretamente dos RPCs: \texttt{POST /v1/rooms}, \texttt{POST /v1/rooms/\{id\}/join},
\texttt{GET /v1/rooms/\{id\}}, \texttt{POST /v1/rooms/\{id\}/start},
\texttt{POST /v1/rooms/\{id\}/leave}, \texttt{POST /v1/rooms/\{id\}/ready} e
\texttt{POST /v1/match/stream}, além de \texttt{GET /healthz}. A referência completa de mensagens
está em \texttt{docs/messages.md}.

\section{Detalhes de implementação}

O backend usa Go 1.25 e \texttt{google.golang.org/grpc}. Cada serviço é um processo
independente em \texttt{services/} com lógica isolada em \texttt{internal/}, o que permite testes
unitários por serviço (\texttt{go test ./...} cobre Gateway, Lobby e Game). O estado é mantido
em memória: o Lobby guarda salas e jogadores com controle de concorrência por \textit{mutex}; o
Game mantém uma partida autoritativa que valida o \textit{input}, descarta entradas repetidas por
\texttt{input\_sequence}, avança um \textit{tick} e recalcula movimento (limitado), abertura de
baús, dano por arma, zona segura que encolhe e ranking ordenado.

O empacotamento usa um \texttt{Dockerfile} por serviço e um
\texttt{deployments/docker-compose.yml} no qual o Gateway só sobe após Game e Lobby ficarem
saudáveis (\texttt{depends\_on} com \texttt{condition: service\_healthy}). Isso torna a demonstração
reprodutível em qualquer máquina com Docker, com um caminho direto para implantação em VPS.

\section{Desafios}

O principal desafio foi consolidar uma base que havia sido revertida do \texttt{master} e
estava espalhada em \textit{branches}, transformando-a em um monorepo coerente com fronteiras
claras entre contrato, serviço e infraestrutura. Decidir que o código gerado ficaria em
\texttt{gen/} e os contratos versionados em \texttt{proto/<svc>/v1} estabilizou os \textit{imports}.
Outro ponto foi a coordenação do grupo de nove pessoas: definimos \textit{ownership} por squad e
um processo de mudança de contrato para evitar quebras. Ficam para as próximas fases conectar
\texttt{Lobby.StartRoom} ao Game, adicionar o transporte WebSocket em tempo real, logs
correlacionados e a prova de carga com 50 jogadores.

\section{Papéis dos alunos}

O grupo se organiza em três squads por fatia funcional (detalhe em \texttt{docs/roles.md}):
\begin{itemize}[noitemsep]
  \item \textbf{Squad A} --- fluxo canônico de início de partida (Lobby + Game).
  \item \textbf{Squad B} --- contrato e esqueleto distribuído (Gateway, \texttt{.proto}, geração).
  \item \textbf{Squad C} --- demo, empacotamento e relatório (Docker Compose, smoke test, SBC).
\end{itemize}
Cada squad é responsável por escrever a seção do relatório correspondente ao que implementou.
Os nomes dos nove alunos e os \textit{owners} de cada squad estão marcados como
\texttt{PLACEHOLDER} e devem ser preenchidos antes da submissão.

\section{Conclusão}

A Entrega 1 cumpre seu objetivo: um esqueleto distribuído executável que demonstra, de forma
verificável, comunicação interna por gRPC e \textit{web services} HTTP, com documentação de
arquitetura e mensagens. A base está pronta para evoluir para o pipeline em tempo real, a
jogabilidade completa e a observabilidade exigida na Entrega 2.

\begin{thebibliography}{9}
\bibitem{grpc} gRPC Authors. \textit{gRPC: A high performance, open source universal RPC framework}. \url{https://grpc.io}.
\bibitem{grpcgateway} gRPC-Gateway Authors. \textit{grpc-gateway: gRPC to JSON proxy generator}. \url{https://github.com/grpc-ecosystem/grpc-gateway}.
\bibitem{protobuf} Google. \textit{Protocol Buffers}. \url{https://protobuf.dev}.
\end{thebibliography}

\end{document}
